% Boost Converter
% This DC-DC converter increases the voltage level from input to output

% Designed by: Amir Ostadrahimi

\documentclass [border=5pt]{standalone}
\usepackage{tikz}
\usepackage[american,cuteinductors,smartlabels]{circuitikz} % A package to draw electrical networks with TikZ

%-- the dimensions of the elements can be changed here
\ctikzset{bipoles/thickness=0.7}
\ctikzset{bipoles/length=1.5cm}
\ctikzset{bipoles/resistor/width=.7}
\ctikzset{bipoles/resistor/height=.25}
\ctikzset{bipoles/diode/height=.3}
\ctikzset{bipoles/diode/width=.3}
\ctikzset{tripoles/thickness=.7}
\ctikzset{bipoles/vsourceam/height/.initial=.7}
\ctikzset{bipoles/vsourceam/width/.initial=.7}
\ctikzset{bipoles/vsourceam/width/.initial=.7}

%settings for fonts and lines
\tikzstyle{every node}=[font=\Large]
\tikzstyle{every path}=[line width=0.9 pt, line cap=round, line join=round]




\begin{document}

 \begin{circuitikz}

%------ Converter
\coordinate  (VBottom) at (0,0); %%VBottom stands for voltage source bottom and is located on (0,0)

\draw (VBottom) to [V, l=$V_{in}$, invert] ++(0,3) coordinate (VTop); % "V" is to insert voltage source, instead of "V", you can use "battery", "battery1", and "battery2".
% the " l=$V_{in}$" is for label. You can use  "l_=$V_{in}$" or "l^=$V_{in}$" to change its location. 
%"invert" changes the polarity of the source. (VTop) stands for voltage top.

\draw (VTop) to [short] ++(1,0) to [L, l=$L$, v=$v_L$, i>^=$i_L$]++ (2,0) coordinate (InductorRight);  % "i>^=$i_L$ " is for showing the current of the element, by using "^" we can change its location and bring it at the input and the output of the element. By using  "<"  and  ">"  we can change its direction. And finally by using "_" we can change its vertical location.

\draw (InductorRight) to [short]++(0,-0.8) node [nigbt, anchor=C, label={ [xshift=-1.2cm, yshift=0] \normalsize $IGBT$}] (nigbt1){}; % here we used an N-channel IGBT using "nigbt". Alternatively, we could use "pigbt", "nmos", "pmos" and so on.

\draw (nigbt1.E)-- (nigbt1.E |-VBottom); % using this command, a line starts from the Emitter of the IGBT and extends vertically until reaches the intersection of the vertical line from (nigbt1.E) and the horizontal line from (VB).

\draw (InductorRight) to [D*] ++(2,0) coordinate (DiodeRight) to [C, l_=$C$, i>^=$i_c$](DiodeRight|-VBottom); % using [D*] makes the diode solid black. If we just use [D], the diode will be white.

\draw 
(DiodeRight) to [short]++(2,0) coordinate(ResistorTop) to [R, l_=$R$, v^=$V_{out}$](ResistorTop|-VBottom) coordinate (ResistorBottom)

(ResistorBottom)-- (VBottom);

%------ Conversion Ratio

\coordinate [label={ [xshift=0, yshift=0] \large $ M(D)=\frac{V_{out}}{V_{in}}=\frac{1}{1-D};$ \large $0 \leq D < 1$ }] (M) at (3,-1);

%------ Curve

\def\xo{2} % Axes of origin
\def\yo{-7} % Axes of origin
\def\length {5} % length of the axes
\def\N{200} %number of samples

\coordinate [label={ [xshift=-5, yshift=-12]  \normalsize$0$ }] (origin) at (\xo,\yo); %

\draw[->] (\xo -0.2, \yo) --++(\length,0) node[right] {\small $D$};
 \draw[->]  (\xo,\yo -0.2) -- ++(0, \length) node[above] {\small $M(D)$};
\draw     plot[samples=\N, domain=  0     :   0.77, xshift=\xo cm, yshift=\yo cm] (5*\x  ,{1/(1-\x)});



\end{circuitikz}



\end{document}